\documentclass[11pt]{article}

\usepackage[margin=1in]{geometry}
\usepackage{amsmath,amssymb,amsthm,mathtools}
\usepackage{hyperref}
\usepackage{enumitem}

\newtheorem{theorem}{Theorem}[section]
\newtheorem{proposition}[theorem]{Proposition}
\newtheorem{lemma}[theorem]{Lemma}
\newtheorem{definition}[theorem]{Definition}
\newtheorem{remark}[theorem]{Remark}

\newcommand{\dR}{\operatorname{dR}}
\newcommand{\Conf}{\operatorname{Conf}}
\newcommand{\Tr}{\operatorname{Tr}}
\newcommand{\Ad}{\operatorname{Ad}}

\title{A Finite\, Spine for Modular Time: de~Rham Monodromy, Radon--Nikodym Structure,
and TKK Grade Preservation}
\author{Nick Goutev (Lean/SymPy formalization team)}
\date{June 2026}

\begin{document}

\maketitle

\begin{abstract}
We present a theorem-honest Lean 4 formalization of the final modular-time bridge
layer of a larger architecture relating non-isotropic configuration-space topology,
operator-algebraic modular dynamics, and thermodynamic clock emergence.

The finite algebraic core proves that the forbidden-cone de~Rham class for the
three-point non-isotropic configuration is identified with the modular derivation
sector compatible with the 5-graded Tits--Kantor--Koecher (TKK) decomposition.
We isolate unresolved analytic ingredients as explicit interface sockets, and include
independent SymPy witness scripts for the key local matrix identities.
\end{abstract}

\section{Introduction}
The emergence of time from state-dependent modular flow (Connes--Rovelli thermal
time) remains a central conceptual issue in quantum statistical mechanics and
noncommutative geometry. Here we provide a finite, computation-first formulation
where the main structural claims are implemented in Lean 4, and analytic completion
is cleanly factorized.

We focus on the final bridge layer:
\[
\text{de~Rham generator around }
\{Q=0\} \,\rightsquigarrow\, \text{Tomita--Takesaki modular derivation}
\rightarrow \text{parabolic clock}.
\]

\section{Topological finite core}
Let $\Conf_3^Q(\mathbb C^D)$ be the 3-point non-isotropic configuration space of an
ordered triple in $\mathbb C^D$ with forbidden light-cone factors $Q(x_i-x_j)=0$.
The formal development proves finite combinatorial invariants in the four-dimensional
case, including generator bookkeeping for the edge formalism and rank targets used in
subsequent modular transport arguments.

\begin{proposition}[Formal finite spine]
In the compiled finite layer, the quadratic non-isotropic de~Rham skeleton and its
cooperadic edge bookkeeping satisfy the expected degree assignments and rank targets
for the targeted branch, including the 32-dimensional formal rank in the calibrated
state.
\end{proposition}

\section{De~Rham form and parabolic contraction}
In coordinates where the light-cone function is
\[
Q = E^2-p_x^2-p_y^2-p_z^2,
\]
we define the logarithmic form $\omega=d\log Q$.  A SymPy witness checks the local
Euler contraction identity
\[
\iota_X \omega = 2,
\]
and vanishing differential of the clock one-form in this affine chart.

\section{Radon--Nikodym and modular derivations}
The central Lean bridge theorem identifies the de~Rham class with the modular
flow generator:
\begin{equation}
\delta(A)=[K,A],
\end{equation}
for the modular Hamiltonian $K$ in the finite presentation.

A set of theorems prove grade invariance of this flow:
\[
\delta \in \mathfrak g_0,\qquad
\delta\,\mathfrak g_i \subset \mathfrak g_i,
\]
and compatibility with the parabolic clock equations used in the final
identification layer.

\section{Parabolic clock and unipotent group law}
The parabolic clock is modeled by a nilpotent shear family
\[ P(t)=I+tN,\quad N^2=0,\]
so that
\[P(t)\circ P(s)=P(t+s),\quad \det P(t)=1.\]
This unipotent, volume-preserving law is verified by symbolic computation and
matches the finite modular-time group law used by the Lean bridge declarations.

\section{Layering and theorem-honesty policy}
The project separates:
\begin{itemize}[leftmargin=*]
  \item\textbf{Lean-honest finite algebra/topology}: de~Rham formal spine,
  cooperad combinatorics, TKK grading checks, and parabolic clock identities.
  \item\textbf{Lean-to-socket transport}: finite-to-analytic comparison maps,
  full Tomita--Takesaki realizability hypotheses, full Gysin/Dupont machinery,
  and physical interpretation hypotheses.
  \item\textbf{Computational audits}: SymPy scripts for local matrix and residue
  identities.
\end{itemize}

\section{Artifacts and reproducibility}
All key artifacts are version-controlled in the repository:
\begin{itemize}[leftmargin=*]
  \item Lean modules:
  `ModularRadonNikodymJacobianBridge`, `ModularTimeDeRhamBridge`,
  `ModularParabolicTimeBridge`.
  \item Graph extraction:
  `proofs/tools/ExtractGraph.lean` and generated `proof_graph.json`.
  \item Symbolic scripts:
  `proofs/modular_time_derivation_witness.py`,
  `proofs/modular_parabolic_time_witness.py`,
  `proofs/modular_radon_nikodym_jacobian_bridge.py`.
\end{itemize}

\section{Conclusion}
This release does not claim a complete analytic realization of all physical
hypotheses. It does provide a formally certified finite core in which the
modular-time identification is expressed as a precise topological--algebraic map.
This is, in particular, a concrete candidate for theorem-honest modular time
authenticity in high-level mathematical physics.

\end{document}
